\documentclass[sigconf]{acmart}

% ACM sigconf packages
\usepackage{booktabs}
\usepackage{array}
\usepackage{url}
\usepackage{graphicx}
\usepackage{xcolor}
\usepackage{enumitem}

% TikZ for diagrams
\usepackage{tikz}
\usepackage{pgfplots}
\pgfplotsset{compat=1.18}
\usetikzlibrary{shapes,arrows,positioning,calc}

% ACM metadata
\setcopyright{acmcopyright}
\copyrightyear{2025}
\acmYear{2025}
\acmDOI{10.1145/XXXXXXX.XXXXXXX}

% Conference information
\acmConference[KDD '25]{31st ACM SIGKDD Conference on Knowledge Discovery and Data Mining}{August 3--7, 2025}{Toronto, ON, Canada}
\acmBooktitle{Proceedings of the 31st ACM SIGKDD Conference on Knowledge Discovery and Data Mining (KDD '25), August 3--7, 2025, Toronto, ON, Canada}

% Bibliography
\bibliographystyle{ACM-Reference-Format}

% CCS concepts
\begin{CCSXML}
<ccs2012>
<concept>
<concept_id>10010147.10010178.10010179</concept_id>
<concept_desc>Computing methodologies~Artificial intelligence</concept_desc>
<concept_significance>500</concept_significance>
</concept>
<concept>
<concept_id>10010147.10010178.10010179.10010182</concept_id>
<concept_desc>Computing methodologies~Multi-agent systems</concept_desc>
<concept_significance>500</concept_significance>
</concept>
</ccs2012>
\end{CCSXML}

\ccsdesc[500]{Computing methodologies~Artificial intelligence}
\ccsdesc[500]{Computing methodologies~Multi-agent systems}

% Keywords
\keywords{Predictive Maintenance, Prognostics and Health Management, Multi-Agent Systems, Benchmark, Agentic AI, Tool Use, ReAct Framework}

\begin{document}

\title{PHM-Agent Benchmark: Expert-Vetted Multi-Asset Predictive Maintenance Evaluation for Agentic Systems}

% Author information (placeholder - to be filled with actual author details)
\author{Research Team}
\affiliation{
  \institution{Institution Name}
  \city{City}
  \country{Country}
}
\email{email@institution.edu}

\begin{abstract}
Predictive maintenance (PHM) in industrial settings requires sophisticated reasoning across multi-modal sensor data, regulatory frameworks, and business constraints. While large language models (LLMs) demonstrate impressive capabilities in general domains, their effectiveness in specialized industrial workflows remains uncertain. We introduce PHM-Agent Benchmark, a comprehensive evaluation framework comprising 75 expert-vetted scenarios spanning five task categories: RUL Prediction, Fault Classification, Engine Health Analysis, Cost-Benefit Analysis, and Safety/Policy Evaluation. Unlike automated scenario generation approaches, our benchmark employs systematic SME-driven curation across 18 datasets covering aircraft engines, industrial bearings, gearboxes, and electric motors. Each scenario reflects authentic operational requests with domain-specific terminology, multi-asset complexity, and empirically-derived ground truth enabling deterministic evaluation. Our systematic scenario generation process---encompassing progressive dataset curation, human-in-the-loop query formulation, and rigorous ground truth extraction---required 396 person-hours of expert effort, ensuring real-world relevance and evaluation rigor. We evaluate both single-agent and multi-agent architectures across multiple LLM backends, establishing performance baselines and identifying critical capabilities for industrial deployment.
\end{abstract}

\maketitle

\section{PHM Benchmark: Scenario Design and Dataset Curation}
\label{sec:scenarios}

This section describes our systematic methodology for creating a comprehensive benchmark of 75 expert-vetted predictive maintenance scenarios. Rather than simply presenting the final scenario collection, we explain the rigorous process we followed: scenario categorization grounded in established PHM literature, progressive dataset curation protocols, SME-driven query formulation pipelines, empirical ground truth extraction, and comprehensive scenario statistics. This process-oriented narrative demonstrates the scientific rigor underlying our benchmark contribution.

%==============================================================================
\subsection{Scenario Categorization: Mapping to PHM Core Tasks}
\label{subsec:categorization}
%==============================================================================

The foundation of our scenario categorization emerges from established Prognostics and Health Management (PHM) literature~\cite{phmBench2024}, which has identified four core task categories that span the asset lifecycle and organizational hierarchy in industrial maintenance operations. These categories---\textit{Condition Monitoring}, \textit{Fault Diagnosis}, \textit{Fault \& RUL Prediction}, and \textit{Maintenance Scheme}---represent fundamentally different types of maintenance decisions that organizations make daily, each requiring distinct analytical capabilities, multi-modal reasoning, and stakeholder-aligned task design. This hierarchical categorization enables systematic evaluation of agentic architectures across the decision-making hierarchy present in industrial operations.

Our five scenario types map systematically to these four established PHM categories through a task taxonomy that respects both technical requirements and organizational structure. \textbf{Engine Health Analysis} (30 scenarios, 40\%) aligns with \textit{Condition Monitoring}, focusing on multi-modal sensor interpretation, anomaly detection, and health status assessment across turbofan engine systems. This category tests agent capabilities in understanding temporal patterns in sensor data (temperature, pressure, vibration), reasoning about operational contexts, and making diagnostic determinations across four cognitive stages: Understanding (signal interpretation), Perception (health classification), Reasoning (trend analysis), and Decision-Making (maintenance recommendations). The dominance of this category reflects the complexity and diversity of condition monitoring tasks in practice, where operators continuously assess equipment health through heterogeneous sensor modalities.

\textbf{Fault Classification} (15 scenarios, 20\%) directly implements \textit{Fault Diagnosis}, requiring agents to identify specific fault types (inner race, outer race, ball damage for bearings; rotor bar failures for motors; component-specific faults for planetary gearboxes) from sensor signatures. This technical core of PHM demands pattern recognition across diverse fault taxonomies, multi-class decision-making under uncertainty, and precise diagnostic capabilities essential for repair planning. Equipment-specific fault localization---determining not just that a failure occurred, but which component failed and how---represents critical diagnostic reasoning that maintenance technicians perform when troubleshooting malfunctioning assets.

\textbf{RUL Prediction} (15 scenarios, 20\%) embodies \textit{Fault \& RUL Prediction}, tasking agents with forecasting remaining operational life measured in cycles, hours, or degradation percentages. This category tests time-series modeling, degradation trend analysis, and quantitative prediction accuracy through error metrics (MAE, RMSE) computed against empirically-derived ground truth. RUL predictions enable proactive maintenance scheduling, spare parts inventory optimization, and operational risk mitigation---capabilities essential for reliability engineers and maintenance planners operating at the site engineering level.

The final two categories---\textbf{Cost-Benefit Analysis} (5 scenarios, 7\%) and \textbf{Safety/Policy Evaluation} (10 scenarios, 13\%)---together constitute the \textit{Maintenance Scheme} category, addressing strategic and compliance considerations that extend beyond technical diagnostics into business logic and regulatory frameworks. Cost-Benefit Analysis integrates RUL predictions with financial modeling (labor overhead rates of 160-200\%, equipment-specific failure costs, downtime impacts) and ROI optimization, serving plant managers, financial analysts, and executives who must balance maintenance investments against operational risks. Safety/Policy Evaluation encompasses regulatory compliance assessment (OSHA vibration limits, FAA airworthiness directives, IEC electrical standards, NEMA motor specifications), risk classification across four levels (low, medium, high, critical), and actionable safety protocol generation. This category tests agent reasoning about proximity hazards, worker safety, equipment safety, and internal company protocols that often exceed regulatory minimums.

This organizational mapping is critical because it reveals \textit{who} within an industrial organization would pose each type of question and how agents must adapt their responses to different stakeholder needs through tool orchestration and multi-hop reasoning. A maintenance technician at the asset management level asking about bearing faults requires immediate, actionable diagnostic information with specific fault characteristics and repair procedures. In contrast, an executive at the management level evaluating cost-benefit tradeoffs needs strategic analysis with comprehensive ROI calculations, risk assessments, and business impact quantification. Our scenario categorization respects these distinctions through stakeholder-aligned task design, ensuring that each scenario reflects the authentic vocabulary, decision priorities, and operational context of its intended user persona. This stakeholder-centric approach enables evaluation of agent capabilities in serving diverse industrial roles rather than optimizing for a single idealized user.

%==============================================================================
\subsection{Dataset Selection and Curation Protocol}
\label{subsec:datasets}
%==============================================================================

Our dataset selection process followed a systematic protocol inspired by MLE-Bench's~\cite{patil2024mlebench} approach of progressive curation from a broad landscape to a focused, high-quality subset. Unlike benchmarks that filter down from massive collections, we adopted an \textit{expansion strategy}: beginning with a single proof-of-concept scenario and progressively growing to 75 scenarios through deliberate multi-asset coverage and task diversity. This expansion narrative demonstrates our commitment to dataset-task alignment, ground truth verification, and multi-modal coverage essential for rigorous agentic evaluation.

We began our curation process with PDMBench's standardized collection of 14 predictive maintenance datasets~\cite{pdmBench2023}, which provided a comprehensive foundation spanning multiple equipment types (bearings, motors, gearboxes, turbofan engines) and operational conditions (varying speeds, loads, temperatures, degradation modes). However, rather than using all available datasets indiscriminately, we applied rigorous dataset-task matching criteria to ensure that each scenario could be properly evaluated with appropriate ground truth. This principle of alignment between task requirements and dataset capabilities guided our entire selection protocol.

\textbf{Expansion Stage 1 (1 scenario):} We initiated with a single proof-of-concept RUL prediction scenario using the CMAPSS (Commercial Modular Aero-Propulsion System Simulation) FD001 dataset~\cite{saxena2008phm}, NASA's turbofan engine degradation simulation. This dataset provides explicit RUL ground truth files containing true remaining cycles for each test unit, enabling quantitative evaluation through Mean Absolute Error (MAE) and Root Mean Squared Error (RMSE) metrics. The initial scenario validated our evaluation framework, ground truth extraction pipeline, and agent orchestration design.

\textbf{Expansion Stage 2 (10 scenarios):} Building on the proof-of-concept, we expanded to 10 scenarios by incorporating FEMTO bearing degradation data (RUL as percentage of lifetime) and CWRU (Case Western Reserve University) bearing fault classification data (inner race, outer race, ball damage with varying severity levels). This expansion introduced multi-asset generalization (engines → bearings) and task diversity (RUL prediction + fault classification), testing agent capabilities across different equipment types and analytical workflows.

\textbf{Expansion Stage 3 (20 scenarios):} The third expansion introduced Cost-Benefit Analysis and Safety/Policy Evaluation scenarios, requiring integration of technical predictions with business logic and regulatory frameworks. We incorporated industry-standard cost models (labor overhead 160-200\%, equipment-specific failure costs) and regulatory thresholds (OSHA, FAA, IEC, NEMA standards) to enable objective evaluation of strategic reasoning. This stage broadened our benchmark beyond technical diagnostics into decision-making domains critical for industrial deployment.

\textbf{Expansion Stage 4 (40 scenarios):} Progressive expansion to 40 scenarios focused on multi-dataset coverage within each equipment class to ensure cross-platform generalization. For bearings, we added XJTU, HUST, MFPT, Mendeley, and Paderborn datasets, each offering distinct characteristics: HUST provides varying speeds and loads, MFPT features ultra-high sampling rates (97.656 KHz) for frequency-domain analysis, Mendeley captures non-stationary conditions with varying RPM (600-1800), and Paderborn offers multimodal sensors (vibration, current, torque, speed). For motors, we incorporated ElectricMotorVibrations (triaxial acceleration, torque, resistance measurements) and RotorBrokenBar (1-4 broken bar severity levels). Gearbox scenarios utilize GearboxUoC and PlanetaryPdM datasets for component-specific fault localization (sun gear, planet gear, ring gear). This diversity ensures agents are evaluated on their ability to generalize across operational contexts, sensor modalities, and equipment architectures rather than overfitting to a single asset class or data distribution.

\textbf{Expansion Stage 5 (75 scenarios):} The final expansion incorporated EngineMTQA~\cite{lei2024itformer}, a large-scale engine health question-answering dataset with over 110,000 expert-validated questions derived from multi-modal sensor data. This addition enabled comprehensive evaluation of Engine Health Analysis scenarios (30 scenarios, 40\% of benchmark) across four cognitive stages: Understanding (open-ended signal interpretation), Perception (closed-ended component health detection), Reasoning (trend prediction and causal analysis), and Decision-Making (maintenance planning and operational recommendations). We also added Azure Predictive Maintenance dataset for additional multi-operational scenario coverage across varying flight conditions and degradation patterns.

The core principle guiding our dataset selection throughout this expansion was \textit{alignment between task requirements and dataset capabilities}. For RUL Prediction scenarios, we exclusively selected datasets containing actual RUL ground truth---specifically, CMAPSS (FD001-FD004) turbofan engine degradation datasets and FEMTO bearing degradation data. These datasets provide timestamped operational cycles or degradation percentages with known failure points, enabling quantitative evaluation of RUL predictions through empirically-derived error metrics. Attempting to evaluate RUL predictions on datasets lacking temporal degradation data with verified ground truth would produce unreliable results and undermine benchmark validity---a critical distinction for deterministic evaluation of agentic systems.

Similarly, for Fault Classification scenarios, we selected datasets explicitly designed for fault diagnosis with labeled fault types and comprehensive fault taxonomies. The CWRU bearing dataset provides multiple bearing fault categories with varying severity levels and operating conditions. The XJTU, HUST, MFPT, Paderborn, and Mendeley bearing datasets offer diverse fault types under varying operational conditions (speed, load, temperature), testing agent robustness to environmental variations. For electric motor faults, ElectricMotorVibrations and RotorBrokenBar datasets capture motor-specific failure modes (rotor bar fractures, bearing degradation, electrical faults). Gearbox faults are represented through GearboxUoC and PlanetaryPdM datasets, enabling component-specific fault localization in planetary gearbox systems through vibration signature analysis and operational parameter monitoring.

This progressive expansion process---from a single proof-of-concept to a curated collection of 75 scenarios across 18 datasets with verified ground truth---ensures that our evaluation framework maintains scientific rigor while achieving comprehensive coverage of diverse industrial equipment types, sensor modalities, operational contexts, and maintenance decision-making workflows. The final dataset distribution spans aircraft engines (CMAPSS FD001-004, Azure, EngineMTQA), industrial bearings (CWRU, FEMTO, IMS, XJTU, HUST, MFPT, Mendeley, Paderborn), gearboxes (GearboxUoC, PlanetaryPdM), and electric motors (ElectricMotorVibrations, RotorBrokenBar), demonstrating multi-asset generalization essential for real-world industrial deployment where maintenance teams manage heterogeneous equipment fleets.

%==============================================================================
\subsection{Scenario Generation Pipeline: SME-Driven Query Formulation}
\label{subsec:generation}
%==============================================================================

The generation of our 75 expert-vetted scenarios followed a structured, SME-driven process that prioritized real-world relevance, domain authenticity, and practical applicability through human-in-the-loop validation. Unlike approaches that rely solely on automated generation or synthetic task creation, our methodology centered on Subject Matter Expert (SME) involvement throughout the scenario formulation process, drawing on expertise from predictive maintenance engineers, operations managers, and domain specialists in aerospace and manufacturing. This approach, inspired by IndustryEQA's human filtering and reannotation methodology~\cite{li2025industryeqa}, ensures that our scenarios test not just technical capabilities but also natural language understanding across different professional vocabularies and operational contexts.

The scenario generation process began with \textit{query formulation}, where SMEs drafted natural language questions representing authentic operational requests they encounter in industrial settings. We established specific guidelines for SMEs to follow during this formulation process, creating a systematic protocol for query development. First, queries should incorporate domain-specific terminology and acronyms as they are actually used in practice---terms like HPC (High-Pressure Compressor), LPC (Low-Pressure Compressor), HPT (High-Pressure Turbine), LPT (Low-Pressure Turbine), MAE (Mean Absolute Error), RMSE (Root Mean Squared Error), RUL (Remaining Useful Life), CMAPSS (Commercial Modular Aero-Propulsion System Simulation), and regulatory standards (FAA, EASA, OSHA, IEC, NEMA, AGMA) appear naturally in industrial discourse. Second, questions should be framed in the voice of real stakeholders, reflecting how plant managers, maintenance technicians, or safety officers would actually pose requests with operational urgency and business context. For example, rather than asking "Calculate the RUL for units in test\_FD001.txt and report MAE against ground truth," an authentic query might ask "Which aircraft engines in our test fleet are running on fumes and need immediate attention before they fail catastrophically?"

This emphasis on authentic domain vocabulary challenges agents to perform query disambiguation and extract intent from informal, context-rich queries that may omit technical details while emphasizing operational urgency through colloquialisms ("running on fumes," "catastrophic failure risk," "immediate intervention required"). Industrial communication is rarely formal or standardized---operators use equipment-specific shortcuts, managers focus on business implications and cost constraints, and technicians employ fault-specific jargon developed through hands-on experience. Our scenarios deliberately capture this linguistic diversity, testing agent capabilities in natural language understanding across heterogeneous professional vocabularies rather than sanitized, academically-formatted prompts.

\textbf{Tool Orchestration and Distraction Tool Testing:} A critical aspect of our scenario design involves evaluating agent capabilities in appropriate tool selection through multi-hop reasoning. Each scenario defines a set of \textit{required tools}---functions the agent must invoke to complete the task successfully (e.g., load\_dataset, train\_rul\_model, predict\_rul, calculate\_mae, verify\_ground\_truth for RUL prediction scenarios). However, to rigorously test tool usage accuracy, scenarios also include \textit{distraction tools}---plausible but inappropriate functions that should not be invoked for the specific task at hand. For instance, a bearing fault classification scenario might include distraction tools such as weather\_data\_loader, financial\_calculator, or web\_search\_tool alongside the correct tools (load\_dataset, train\_fault\_classifier, classify\_faults, calculate\_accuracy). This design tests whether agents can discriminate between contextually appropriate and inappropriate tools, a critical capability for industrial deployment where systems have access to extensive function libraries but must select only relevant operations. Tool usage penalties in our evaluation framework (Section 5) explicitly penalize invocation of distraction tools, rewarding agents that demonstrate precise tool orchestration and reasoning about task-tool alignment.

Our scenarios exhibit significant structural diversity across multiple dimensions that reflect real-world information needs. \textbf{Query types} range from open-ended questions requiring analytical synthesis ("What factors are contributing to premature bearing failures in high-load applications?") to close-ended multiple-choice questions with deterministic answers ("Does this equipment meet OSHA vibration standards---yes or no?"). This distinction mirrors operational reality: sometimes stakeholders need exploratory analysis and diagnostic reasoning, while other times they need definitive compliance determinations or binary health assessments. We intentionally include both types (60\% open-ended, 40\% close-ended) to evaluate agent reasoning capabilities across different response modalities and task structures.

Another critical dimension is \textbf{data provision mode}, which tests agent capabilities in analytical planning and dataset navigation. Some scenarios embed all necessary data directly within the query itself, providing sensor readings, operational parameters, or maintenance histories as part of the question. For example, an EngineMTQA scenario might present a complete time-series of temperature, pressure, and vibration measurements and ask "Based on these sensor readings from turbofan unit 47, what is the engine health status?" This \textit{embedded data mode} (45\% of scenarios) focuses evaluation on analytical reasoning, sensor interpretation, and diagnostic capabilities given complete information. Other scenarios deliberately omit specific data references, requiring agents to determine which dataset to load (e.g., CMAPSS\_FD001 vs. FEMTO\_bearing\_1), which features to extract (temporal vs. frequency-domain), and which analysis to perform. This \textit{data discovery mode} (55\% of scenarios) tests agent capabilities in tool selection, dataset navigation, and analytical planning---critical for real-world deployment where queries often don't specify all operational details. The distinction represents two common industrial personas: domain experts who provide complete context, and operational managers who pose high-level questions expecting agents to gather necessary data autonomously.

The \textbf{multi-asset versus single-asset} distinction introduces another layer of complexity testing agent capabilities in prioritization and fleet-level reasoning. Single-asset scenarios focus on analyzing one piece of equipment in isolation, such as predicting the RUL of a specific bearing (unit 14 in FEMTO dataset) or diagnosing a fault in a particular motor (rotor bar configuration in ElectricMotorVibrations). Multi-asset scenarios (30\% of benchmark) present fleet-level questions where agents must first identify which assets within a larger population require attention. For instance, "Which of our 100 turbofan engines will fail in the next 30 days, and which 10 should receive immediate maintenance?" requires the agent to process an entire fleet, compute RUL predictions for all units, rank by risk (lowest RUL = highest priority), apply operational threshold criteria (30-day window), and identify the critical subset requiring immediate intervention. This mirrors real operational scenarios where maintenance teams manage portfolios of heterogeneous assets with limited resources and must prioritize effectively through multi-asset generalization and risk-based ranking.

The complete scenario generation process involved approximately 396 person-hours across three primary roles, demonstrating substantial human-intensive investment in benchmark quality. A primary scenario writer with expertise in predictive maintenance and data science spent approximately 200 hours drafting technical scenarios, selecting appropriate datasets through systematic dataset-task alignment, and defining ground truth validation criteria with empirical calibration. Two SME reviewers---one representing operations and maintenance leadership (100 hours), the other bringing aerospace and manufacturing domain expertise (96 hours)---contributed to validating real-world relevance, ensuring technical accuracy, and refining scenarios through multiple review iterations with expert annotation. This labor-intensive process, totaling \$63,360 in expert effort at \$160/hour overhead rates, ensures that every scenario reflects genuine operational needs, uses authentic domain vocabulary with appropriate technical terminology density, and can be objectively evaluated against verified ground truth through deterministic evaluation protocols rather than artificial or contrived tasks optimized for LLM performance.

%==============================================================================
\subsection{Ground Truth Utilization and Validation Framework}
\label{subsec:groundtruth}
%==============================================================================

Establishing rigorous ground truth for each scenario was essential to enable deterministic, reproducible evaluation of agentic systems through empirical calibration. Unlike open-ended generative tasks where evaluation may be subjective, our scenarios require verifiable correctness criteria that can be automatically assessed through threshold-based evaluation against validated references. The ground truth preparation and utilization process varied significantly across scenario types, reflecting the different evaluation requirements for RUL prediction (quantitative error metrics), fault classification (categorical accuracy), cost-benefit analysis (financial modeling), safety evaluation (regulatory compliance), and engine health analysis (expert-validated responses).

\textbf{RUL Prediction Ground Truth Extraction:} For RUL Prediction scenarios, we extracted actual RUL values directly from the datasets themselves through systematic parsing protocols. The CMAPSS datasets (FD001-FD004) provide explicit RUL files (RUL\_FD001.txt, RUL\_FD002.txt, RUL\_FD003.txt, RUL\_FD004.txt) containing the true remaining cycles for each test unit at the final operational timestamp. We parsed these files to establish expected MAE and RMSE ranges based on actual data distributions through empirical analysis. For example, for CMAPSS\_FD001, we determined through statistical analysis of ground truth values that MAE values between 12-20 cycles and RMSE values between 18-26 cycles represent realistic prediction performance for state-of-the-art models, with mean RUL of 75.5 cycles and range spanning 7-145 cycles across test units. The FEMTO dataset required different handling, as it provides RUL as a percentage of lifetime degradation rather than absolute cycles. We developed extraction logic to identify unique bearing segment IDs and compute expected RUL percentage distributions, establishing validation thresholds of 5-15\% MAE for bearing degradation prediction based on empirical calibration against actual test trajectories.

This emphasis on extracting real ground truth through deterministic evaluation protocols rather than estimating expected values addresses a critical limitation in prior benchmark work. Initial versions of our scenarios contained placeholder ranges like "MAE should be approximately 15-20 cycles" without verification against actual data distributions. We systematically replaced all such estimates with empirically derived values, ensuring that validation criteria reflect genuine dataset characteristics through rigorous ground truth verification rather than assumptions or heuristic approximations. This methodological rigor enables confident assessment of agent prediction accuracy against validated references.

\textbf{Fault Classification Ground Truth:} For Fault Classification scenarios, ground truth preparation involved validating fault labels across diverse datasets with varying fault taxonomies and severity categorizations. The CWRU dataset categorizes bearing faults into inner race defects, outer race defects, and ball element damage with multiple severity levels (0.007", 0.014", 0.021" crack depths) and operating conditions (motor loads, shaft speeds). We established expected accuracy ranges (typically 85-95\%) based on published benchmarks for these datasets, reflecting state-of-the-art performance on standard fault classification tasks. The ElectricMotorVibrations dataset required special handling because it lacks a segment\_idx column present in other datasets; we adapted our extraction logic through conditional processing to handle both dataset structures, ensuring broad applicability across heterogeneous data formats. The RotorBrokenBar dataset provides a distinct fault taxonomy (healthy baseline, 1 broken bar, 2 broken bars, 3 broken bars, 4 broken bars) with increasing severity levels, requiring custom validation logic for motor-specific fault types and multiclass classification across discrete damage states.

\textbf{Cost-Benefit and Safety Ground Truth:} Cost-Benefit Analysis and Safety/Policy Evaluation scenarios posed unique challenges because ground truth often involves industry standards, regulatory thresholds, or financial models rather than dataset-derived values. For cost-benefit scenarios, we established expected cost ranges based on typical maintenance costs for different equipment types and failure modes, calibrated to industry-standard labor overhead rates (160-200\% rather than unrealistic 50\% rates), equipment-specific failure costs (bearing replacement \$2k-\$5k, motor overhaul \$10k-\$25k, turbofan engine replacement \$500k+), and downtime impact modeling. Safety scenarios reference specific regulatory standards with defined threshold values: OSHA permissible exposure limits (PEL) for vibration, FAA airworthiness directives for cycle limits, IEC 60034 specifications for rotating electrical machines, NEMA MG1 motor standards, AGMA 2001 gear standards. While these scenarios may not have single "correct" answers in the same way RUL predictions do (where ground truth is a specific numerical value), we defined acceptable ranges and compliance criteria that enable objective evaluation through regulatory threshold checking and cost model verification.

\textbf{Engine Health Analysis Ground Truth:} The EngineMTQA scenarios presented a different ground truth challenge due to their multi-modal nature, combining time-series sensor data (temperature, pressure, vibration, RPM) with natural language questions spanning four cognitive stages. Ground truth for these scenarios comes from expert-validated answers in the EngineMTQA dataset itself~\cite{lei2024itformer}, which provides reference responses for questions across Understanding (open-ended signal interpretation: "What does increasing HPC outlet temperature indicate?"), Perception (closed-ended health classification: "Is fan health status normal or fault?"), Reasoning (trend analysis: "Based on LPT cooling trends, what failure mode is developing?"), and Decision-Making (maintenance planning: "Should this engine undergo immediate inspection or continue operation?"). We preserved these validated answers as ground truth while adapting the query format to fit our agentic evaluation framework, ensuring consistency with SME-validated diagnostic reasoning.

\textbf{Ground Truth Utilization in Validation:} Beyond merely establishing ground truth values, our evaluation framework explicitly enforces their utilization through mandatory verification requirements embedded in scenario specifications. Each scenario defines a structured output template that agents must populate, including not only final answers (predicted RUL values, classified fault types, cost analyses, compliance determinations) but also explicit ground truth verification components. For RUL prediction scenarios, agents must compare predictions against known ground truth, calculate error metrics (MAE/RMSE), and validate that computed errors fall within acceptable ranges (\textpm20\% of empirically-derived thresholds). This requirement addresses a common failure mode where agents produce predictions without validation, potentially reporting unrealistic or incorrect results with false confidence.

Our evaluation framework enforces this ground truth utilization requirement through success criteria that demand both task completeness (\textgreater80\% of required components addressed) and successful ground truth verification. A task that produces RUL predictions without validating them against ground truth files is scored as incomplete, even if the predictions themselves happen to be accurate. This design choice enforces best practices in predictive maintenance workflows, where validation against known outcomes is essential to build confidence in model reliability and diagnostic accuracy. The verification requirement ensures that agents not only perform analytical tasks but also validate their outputs against empirical references---a critical capability for industrial deployment where unvalidated predictions pose operational risks.

Validation criteria vary by scenario type but generally assess three dimensions through our deterministic evaluation protocol: \textit{completeness} (were all required task components addressed, including data loading, model training/loading, prediction/classification, metric computation, and verification?), \textit{correctness} (do the answers match or fall within acceptable ranges of ground truth values, with quantitative thresholds for RUL errors and accuracy requirements for classification?), and \textit{efficiency} (was the task completed within reasonable resource constraints without redundant operations or inappropriate tool usage?). This comprehensive ground truth preparation and utilization process ensures that every scenario in our benchmark can be evaluated objectively and consistently through threshold-based evaluation. When an agent claims an MAE of 15 cycles for CMAPSS\_FD001 RUL prediction, we can verify this claim against actual test data through deterministic comparison. When an agent classifies a bearing fault as "inner race damage at 0.014-inch crack depth," we can check this against labeled ground truth through categorical matching. This deterministic evaluation capability distinguishes our benchmark from more subjective evaluation frameworks and enables rigorous comparison across different agent architectures, LLM backends, and tool orchestration strategies.

%==============================================================================
\subsection{Scenario Statistics and Distribution Analysis}
\label{subsec:statistics}
%==============================================================================

Our final benchmark comprises 75 expert-vetted scenarios distributed across 18 datasets, five task categories, and multiple query types, representing a comprehensive evaluation framework for agentic PHM systems. This distribution reflects both the relative importance and complexity of different task types in industrial practice, as well as our systematic expansion strategy from initial proof-of-concept to comprehensive multi-asset, multi-task coverage. Table~\ref{tab:scenario_breakdown} presents the distribution of scenarios across task categories, revealing deliberate choices in scenario allocation driven by operational significance and cognitive complexity.

\textbf{Task Category Distribution and Justification:} Engine Health Analysis scenarios constitute the largest category (30 scenarios, 40\%) because they encompass diverse cognitive sub-tasks spanning four stages of diagnostic reasoning. These scenarios leverage the EngineMTQA dataset's extensive coverage of multi-modal sensor data (temperature, pressure, vibration, RPM across turbofan components: HPC, LPC, HPT, LPT, fan) and expert-validated question-answer pairs. The 30 scenarios distribute across Understanding (16 scenarios testing open-ended signal interpretation and temporal pattern recognition), Perception (8 scenarios testing closed-ended health status classification), Reasoning (4 scenarios testing causal analysis and degradation trend prediction), and Decision-Making (2 scenarios testing maintenance planning and operational recommendations). This cognitive diversity within a single task category justifies its 40\% allocation---it effectively encompasses multiple evaluation dimensions within condition monitoring workflows.

RUL Prediction and Fault Classification receive equal emphasis (15 scenarios each, 20\% apiece), representing the technical core of PHM through multi-task learning and cross-platform evaluation. Despite their equal scenario counts, these categories differ significantly in computational complexity and evaluation rigor through their distinct analytical requirements. RUL scenarios demand time-series modeling across operational trajectories spanning tens to hundreds of cycles, degradation trend analysis through temporal feature extraction, and quantitative error metric computation (MAE/RMSE) against empirically-derived thresholds. Fault classification scenarios emphasize pattern recognition across heterogeneous fault taxonomies (bearing races, motor components, gearbox elements), multi-class decision-making under uncertainty (distinguishing inner race vs. outer race vs. ball damage), and accuracy assessment across varying operational conditions (speeds from 600-1800 RPM, loads, temperatures). The equal allocation reflects their complementary importance: RUL prediction enables proactive planning, while fault classification enables reactive diagnosis---both essential for comprehensive maintenance decision support.

Cost-Benefit Analysis (5 scenarios, 7\%) and Safety/Policy Evaluation (10 scenarios, 13\%) represent smaller but critical scenario categories that test agent capabilities in strategic reasoning and regulatory compliance. While these categories have fewer scenarios, they often require more complex multi-hop reasoning that integrates technical analysis with business logic, financial modeling, regulatory knowledge, and risk assessment frameworks. A single cost-benefit scenario may require RUL predictions for fleet prioritization, failure cost estimation across equipment types (bearing \$2k-\$5k, motor \$10k-\$25k, turbofan \$500k+), maintenance cost modeling with labor overhead (160-200\%), downtime impact quantification, and ROI optimization across preventive versus reactive strategies---synthesizing multiple analytical components into strategic recommendations. Similarly, safety scenarios demand regulatory threshold checking across multiple frameworks (OSHA PEL for vibration, FAA AD for airworthiness, IEC 60034 for electrical systems), four-level risk classification (low, medium, high, critical rather than standard three-level taxonomies), proximity hazard assessment, and actionable safety protocol generation. The smaller allocation (7\% and 13\%) reflects not diminished importance but rather the fact that each scenario integrates extensive domain knowledge and multi-step reasoning, making them individually more comprehensive than single-objective technical tasks.

\textbf{Dataset Distribution and Multi-Asset Coverage:} Table~\ref{tab:dataset_distribution} presents the distribution of scenarios across our 18 curated datasets, organized by equipment type to demonstrate systematic multi-asset generalization. The distribution demonstrates comprehensive coverage across four equipment classes: bearings (32 scenarios across 8 datasets: CWRU, FEMTO, IMS, XJTU, HUST, MFPT, Mendeley, Paderborn), electric motors (8 scenarios across 2 datasets: ElectricMotorVibrations, RotorBrokenBar), gearboxes (5 scenarios across 2 datasets: GearboxUoC, PlanetaryPdM), and aircraft engines (30 scenarios across 6 datasets: CMAPSS FD001-004, Azure, EngineMTQA). This diversity ensures that agents are evaluated on their ability to generalize across equipment architectures, sensor modalities (vibration, current, temperature, pressure, RPM, torque), and operational contexts (varying speeds, loads, environmental conditions) rather than overfitting to a single asset class or data distribution.

The bearing dataset diversity merits particular emphasis, as it represents the most extensive multi-dataset coverage within a single equipment class. The eight bearing datasets offer complementary characteristics for comprehensive evaluation: CWRU provides standard fault types with controlled severity levels, FEMTO enables run-to-failure degradation analysis, IMS offers high-temporal-resolution vibration signatures (20kHz sampling), XJTU captures five distinct fault modes, HUST tests robustness across varying speeds (600-1500 RPM) and loads, MFPT enables frequency-domain analysis through ultra-high sampling rates (97.656 KHz), Mendeley tests non-stationary signal processing with varying speeds (600-1800 RPM), and Paderborn provides multimodal sensor fusion (vibration + current + torque + speed). This systematic diversity within bearings alone tests agent capabilities in cross-platform generalization---a critical requirement for industrial deployment where maintenance teams manage heterogeneous bearing types across different manufacturers, operational settings, and degradation modes.

\textbf{Query Characteristics and Structural Diversity:} Query characteristic statistics further reveal scenario diversity across dimensions that mirror real operational workflows. Of the 75 scenarios, approximately 60\% (45 scenarios) are open-ended analytical questions requiring synthesized explanations ("What factors are contributing to premature bearing failures under high-load conditions in the HUST dataset?"), while 40\% (30 scenarios) are close-ended questions with deterministic answers or multiple-choice options ("Does turbofan unit 23 meet FAA airworthiness requirements for continued operation---yes or no?"). This distribution reflects authentic information needs: sometimes stakeholders need exploratory analysis and diagnostic reasoning, while other times they need definitive determinations for compliance or operational decisions.

Data provision mode statistics demonstrate balanced coverage of two common industrial personas. Approximately 55\% (41 scenarios) require agents to perform data discovery and analytical planning by identifying appropriate datasets (e.g., determining whether to use CMAPSS\_FD001 or FEMTO\_bearing\_1 based on query context), loading relevant data files, and selecting analytical approaches. The remaining 45\% (34 scenarios) embed necessary data within the query itself through sensor readings, operational parameters, or maintenance histories, focusing evaluation on diagnostic reasoning given complete information. This distinction represents operational manager personas (who pose high-level questions expecting autonomous data gathering) versus domain expert personas (who provide complete technical context in their queries).

Multi-asset fleet analysis comprises roughly 30\% (23 scenarios) of the benchmark, requiring agents to process multiple equipment units simultaneously, compute predictions or classifications across the fleet, rank by risk or priority (typically lowest RUL = highest risk), apply operational threshold criteria (e.g., "identify units failing within 30 days"), and generate prioritized maintenance lists. The remaining 70\% (52 scenarios) focus on single-asset analysis with deeper diagnostic detail on individual equipment instances. This 30/70 split reflects actual industrial workflows where fleet management occurs periodically (daily or weekly maintenance planning sessions) while diagnostic troubleshooting occurs continuously (responding to individual equipment alarms or anomalies).

\textbf{Organizational Stakeholder Mapping:} Organizational level mapping reveals that scenarios span all hierarchical levels of industrial decision-making, enabling evaluation of agent adaptability across diverse stakeholder needs. Approximately 35\% of scenarios (26 scenarios) target site engineering and asset management levels, serving maintenance technicians (fault diagnosis with repair procedure recommendations), reliability engineers (RUL prediction for equipment prioritization), and condition monitoring specialists (sensor signal interpretation and anomaly detection). Approximately 40\% (30 scenarios) target management and executive levels, serving plant managers (strategic maintenance planning and resource allocation), financial analysts (cost-benefit optimization and ROI justification), department directors (budget forecasting and capital expenditure planning), and vice presidents (portfolio-level risk assessment). The remaining 25\% (19 scenarios) target regulatory and safety levels, serving environmental health \& safety officers (compliance determination and safety protocol generation), compliance managers (multi-framework regulatory adherence), and risk assessors (four-level risk classification and mitigation strategy development).

This stakeholder distribution ensures that our benchmark evaluates agent capabilities in serving diverse industrial personas with heterogeneous priorities, vocabularies, and decision contexts. A maintenance technician needs immediate, actionable diagnostic information with component-level fault localization ("inner race damage at 0.014-inch crack depth, recommend bearing replacement within 48 hours"). A plant manager needs strategic analysis with business impact quantification ("preventive maintenance on 12 high-risk units costs \$87k but avoids \$340k in reactive failure costs and 156 hours of unplanned downtime"). A safety officer needs regulatory compliance verification with standard citations ("Unit 47 vibration exceeds OSHA PEL of 0.5g RMS, requires immediate operational restriction per 29 CFR 1910.95"). Our scenario design captures these distinct stakeholder-aligned task requirements through authentic vocabulary, appropriate analytical depth, and decision-relevant outputs.

\textbf{Benchmark Evolution and Quality Investment:} The complete scenario generation process, from initial proof-of-concept (1 scenario) through systematic expansion to comprehensive coverage (75 scenarios), represents approximately 396 person-hours of expert effort distributed across three specialized roles. This substantial human-intensive investment demonstrates our commitment to benchmark quality, real-world relevance, and evaluation rigor through progressive curation and expert annotation. A primary scenario writer with dual expertise in predictive maintenance and data science invested approximately 200 hours drafting technical scenarios, performing systematic dataset-task alignment, and defining ground truth validation criteria through empirical calibration against actual data distributions. Two SME reviewers contributed 196 hours combined (100 hours from an operations and maintenance leadership SME, 96 hours from an aerospace and manufacturing domain expert SME), validating real-world operational relevance, ensuring technical accuracy across diverse equipment types, and refining scenarios through multiple review iterations with detailed feedback.

This labor-intensive methodology, totaling \$63,360 in expert effort at industry-standard overhead rates (\$160/hour including benefits, facilities, management), ensures that every scenario reflects genuine operational needs encountered daily in industrial maintenance operations, uses authentic domain vocabulary with appropriate technical terminology density (averaging 4.2 technical terms per query across PHM acronyms, regulatory standards, equipment components, and diagnostic metrics), and can be objectively evaluated against verified ground truth through deterministic protocols rather than subjective assessment. While automated scenario generation methods could produce larger benchmarks more rapidly through LLM-based synthesis, our human-centric approach guarantees that each scenario represents actual industrial workflows, incorporates domain expertise accumulated through years of maintenance experience, and maintains evaluation validity through empirical ground truth rather than synthetic approximations.

The systematic expansion from 1 to 75 scenarios over a six-month development period (January to July 2025) demonstrates our deliberate growth strategy: validate evaluation framework through proof-of-concept (1 scenario, Week 1-2) → expand task diversity (10 scenarios, Week 3-6) → broaden strategic reasoning (20 scenarios, Week 7-12) → achieve multi-asset coverage (40 scenarios, Week 13-20) → incorporate multi-modal cognitive analysis (75 scenarios, Week 21-26). This progressive approach enabled iterative refinement of ground truth extraction protocols, query formulation guidelines, and evaluation criteria based on early validation results, ensuring benchmark quality scaled with coverage expansion.

% Placeholder for tables and figures
% Table 4.1: Scenario Distribution by Task Category [to be added]
% Table 4.2: Dataset Distribution by Equipment Type [to be added]
% Table 4.3: Query Characteristics Breakdown [to be added]
% Table 4.4: Person-Hours Investment by Role [to be added]
% Figure 4.1: Donut Chart - PHM Category Mapping [to be added]
% Figure 4.2: Organizational Hierarchy Diagram [to be added]
% Figure 4.3: Reverse Funnel - Scenario Expansion (1→75) [to be added]
% Figure 4.4: Scenario Generation Pipeline Flowchart [to be added]

\bibliography{references}

\end{document}
